\subsection{Evaluation framework}
\label{sec:method-eval}

\subsubsection{Predicted $\Delta\mathrm{X_{CO2}^{model}}$ evaluation}
\label{sec:method-eval-del-xco2}

On the test subsets of the training/analysis data, we evaluate each model by how
well its predicted anomaly $\Delta\mathrm{X_{CO2}^{model}}$ reproduces the target
$\Delta\mathrm{X_{CO2}^{B11}}$. We report the date-blocked coefficient of
determination ($R^2$) and root-mean-square error (RMSE) of the predicted versus
target anomaly, together with the per-footprint RMSE and residual scatter. Because
the deep ensemble is probabilistic, we also assess the calibration and the
prediction-interval coverage of its intervals. Statistical significance is
established with paired Wilcoxon signed-rank tests and a site-clustered bootstrap,
and we additionally report a random-effects residual comparison and the
associated representation error.

We evaluate all models under a common protocol: the Ridge and XGBoost baselines,
the TabM and deep-ensemble models, the feature-free within-orbit smoother (a
predictor-free null model), and the feature-set ablations
(Sect.~\ref{sec:...}).


\subsubsection{$\mathrm{X_{CO2}}$ comparison across different products}
\label{sec:method-eval-obs-xco2}

To compare OCO-2 $\mathrm{X_{CO2}}$ ($\mathrm{X_{CO2}^{OCO}}$) with $\mathrm{X_{CO2}}$
from TCCON stations ($\mathrm{X_{CO2}^{TCCON}}$), the ATom airborne campaign
($\mathrm{X_{CO2}^{ATom}}$), and two shipborne campaigns
($\mathrm{X_{CO2}^{ship}}$), we apply product-specific harmonization to reduce the
differences that arise from the different retrieval methods. For the TCCON
comparison, we apply the averaging-kernel (AK) and prior harmonization together
with a wet-to-dry prior conversion, following
\citet{wunch2011tccon_network,das2025oco2_v111_tccon_ggg2020,rodgers2003intercomparison}.
The details of the AK and prior harmonization and of the wet-to-dry conversion are
given in Appendix~\ref{app:ak_operation}.

For the ATom comparison, we build a pseudo-column dry-air $\mathrm{X_{CO2}^{ATom}}$
from the in-situ \ce{CO2} profiles measured by the aircraft.

For the shipborne measurements, the averaging kernel and prior they used are not
available, so we compare $\mathrm{X_{CO2}^{OCO}}$ with $\mathrm{X_{CO2}^{ship}}$
directly. The residual product-level difference may therefore be larger for the
shipborne comparison.

We set the collocation search radius between the TCCON station, the ATom
pseudo-column center, or the ship-track center and the OCO-2 soundings to
$100$~km, and the temporal search window to $\pm1$~h around the OCO-2 overpass.
The $\pm1$~h window is the same as \citet{das2025oco2_v111_tccon_ggg2020}, while
our spatial window is stricter than their default coincidence box of
$2.5^\circ \times 5^\circ$ in latitude--longitude centred on the station
($\pm1.25^\circ$ latitude, corresponding to about $139$~km, and $\pm2.5^\circ$
longitude); they tighten the box only at a few sites affected by local pollution
or terrain. Unlike their analysis, which uses only good-quality soundings
(\texttt{xco2\_quality\_flag}~$=0$) and compares the median of each overpass, we
include both quality flags (0 and 1) and evaluate the comparison at the
individual footprint level, which is necessary to resolve the near-cloud
behavior that is the subject of this study.

For each station-day $s$ (one TCCON site on one overpass date) with $n_s$
collocated soundings, let
$e_{s,i} = X_{\mathrm{CO2},i} - X^{\mathrm{TCCON}}_{\mathrm{CO2},i}$
denote the footprint residual against the AK-harmonised TCCON reference,
evaluated once for the uncorrected product and once after correction. The
station-day mean bias and per-footprint root-mean-square error are
\begin{equation}
  b_s \;=\; \frac{1}{n_s}\sum_{i=1}^{n_s} e_{s,i},
  \qquad
  R_s \;=\; \Bigl(\frac{1}{n_s}\sum_{i=1}^{n_s} e_{s,i}^{2}\Bigr)^{1/2}.
\end{equation}
Figure~8a (Results 4.4) summarises the $S = 75$ station-days by three aggregates:
the station-day-equal mean absolute bias
$\overline{|b|} = S^{-1}\sum_s |b_s|$, which weights every station-day
equally regardless of its footprint count; the root-mean-square bias
$\mathrm{RMS}_b = (S^{-1}\sum_s b_s^{2})^{1/2}$, a quadratic aggregate of
the same station-day mean biases that emphasizes the largest-bias
station-days (it is not a footprint-level error metric); and the mean
per-footprint RMSE $\overline{R} = S^{-1}\sum_s R_s$, which measures
footprint-level scatter about the reference. Each row of Fig.~8a shows
the change $\Delta$ (corrected $-$ uncorrected) of one aggregate, so
negative values are improvements: $\overline{|b|}$ and $\mathrm{RMS}_b$
respond only to systematic station-day offsets, whereas $\overline{R}$
also responds to within-overpass scatter. Confidence intervals and
$p$-values come from a site-clustered bootstrap: the 18 sites are
resampled with replacement, each drawn site contributing all of its
station-days ($B = 10\,000$ resamples), which respects the strong
within-site clustering of the sample (e.g.\ R\'eunion contributes 14
station-days); the interval is the 2.5--97.5 percentile range of the
resampled $\Delta$, and the two-sided
$p = 2\min[P(\Delta \ge 0),\, P(\Delta \le 0)]$, floored at
$2/(B+1)$. The paired Wilcoxon test quoted beneath the panel treats
station-days as exchangeable pairs and is reported as a
clustering-agnostic cross-check.